\documentclass[11pt]{article}
\usepackage[a4paper,margin=18mm]{geometry}
\usepackage{fontspec}
\setmainfont{Latin Modern Roman}
\usepackage{amsmath,amssymb,amsthm,amscd,graphicx,color}
\usepackage{xurl}
\usepackage[unicode,hidelinks]{hyperref}
\allowdisplaybreaks
\setlength\emergencystretch{3em}
\numberwithin{equation}{section}

\newtheorem{Theorem}{Theorem}[section]
\newtheorem{Corollary}[Theorem]{Corollary}
\newtheorem{Lemma}[Theorem]{Lemma}
\newtheorem{Proposition}[Theorem]{Proposition}

\def\1{^{-1}}

\def\a{\mathfrak{a}}
\def\al{\alpha}

\def\be{\beta}
\def\bi{{\bar\imath }}
\def\bj{{\bar\jmath }}
\def\bk{{\bar k }}
\def\bl{{\bar l }}
\def\bp{{\bar p }}
\def\bq{{\bar q }}
\def\br{{\bar r }}
\def\bs{{\bar s }}

\def\CC{{\mathbb{C}}}
\def\CMN{\CC^{M|N}}

\def\de{\delta}
\def\De{\Delta}

\def\End{\operatorname{End}}
\def\ep{\varepsilon}

\def\g{\glMN [u]}
\def\ge{\geqslant}
\def\glMN{\mathfrak{gl}_{M|N}}
\def\gr{{\rm gr} }
\def\grpr{{\rm gr}^{\prime} }

\def\id{{\mathrm{id}}}

\def\le{\leqslant}

\def\ns{{\hskip-.5pt}}

\def\om{\omega}
\def\ot{\otimes}

\def\RS{R^{\sharp}}

\def\S{\operatorname{S}}
\def\si{\sigma}
\def\str{\operatorname{str}}
\def\Sym{\mathfrak{S}}

\def\TD{\dot{T}}
\def\TP{T^{\prime}}
\def\TS{T^{\sharp}}

\def\Ua{\operatorname{U}(\a [u])}
\def\UMN{\operatorname{U}(\glMN)}
\def\Ug{\operatorname{U}(\g)}
\def\XP{X^{\prime}}
\def\YMN{\operatorname{Y}(\glMN)}
\def\ZZ{\mathbb{Z}}

\makeatletter\newlabel{ss}{{3.2}{}{Original source locator}{}{}}\makeatother
\begin{document}
\begin{center}\Large For M,N>=1 over C, the degree-(r-1) associated graded RTT Yangian of the general linear Lie superalgebra on an M-even, N-odd superspace is the enveloping algebra of its polynomial current Lie superalgebra\end{center}
\noindent\textbf{Stable ID:} 1651\par
\noindent\textbf{Authors:} Maxim Nazarov\par
\noindent\textbf{Original work:} Yangian of the General Linear Lie Superalgebra\par
\noindent\textbf{Version:} Official SIGMA 16 (2020) article 112, DOI 10.3842/SIGMA.2020.112; journal PDF and native TeX acquired 2026-09-30, exact bytes pinned by SHA256\par
\noindent\textbf{Selected task scope:} Exactly the associative Z/2-graded algebra part of Theorem2.8, with integers M,N>=1 and coefficients C. Give T\_ij\textasciicircum{}(r) filtration degree r-1 and write Y\_ij\textasciicircum{}(r) for its leading class; the isomorphism U(gl(M|N)[u])->gr-prime Y(gl(M|N)) sends e\_ji*u\textasciicircum{}r to -(-1)\textasciicircum{}(parity(j))*Y\_ij\textasciicircum{}(r+1). The RTT/super-sign conventions are those of (2.1)–(2.8). No centre classification, first-filtration freeness as a separate unit, or new Hopf-structure verification is counted.\par
\noindent\textbf{Difficulty:} H2: The substantive obligation is injectivity of the current-algebra map, which is not obtained by merely taking leading terms in the RTT relations. The author separates every possible relation using tensor evaluation representations, supersymmetrization and independent spectral-parameter polynomials, including the identity-current factors that ordinary matrix tensors cannot directly distinguish. This proves a genuine super-Yangian PBW deformation statement beyond the usual qualifying-exam PBW theorem for finite-dimensional Lie algebras.\par
\noindent\textbf{Editorial boundary note (not author text):} The original Theorem 2.8 is retained verbatim. The selected task is its graded associative-algebra isomorphism, established by the complete leading RTT-map calculation and tensor-evaluation separation argument of Propositions 2.6–2.7. All parity, tensor, current, filtration and evaluation conventions needed for that argument are included. The Hopf formulation of the original theorem is retained as source context, not a separate selected proof obligation; centre/Berezinian results and the neighboring first-filtration corollary are excluded.\par
\noindent\textbf{Source:} \url{https://sigma-journal.com/2020/112/}\quad DOI: \url{https://doi.org/10.3842/SIGMA.2020.112}\par
\noindent\textbf{License and attribution:} Copyright retained by the named authors. Published by SIGMA under CC BY 4.0: \url{https://creativecommons.org/licenses/by/4.0/}. Source: \url{https://sigma-journal.com/about.html}. Adaptation consists of standalone excerpt formatting, cover and numbering restoration; original author language and mathematical text are retained.\par
\noindent\textbf{External original locator ss:} Original Proposition 3.2 computes the square of the antipode. It is mentioned only in the unselected Hopf-structure context; the selected graded associative-algebra proof uses Propositions 2.6–2.7, included in full.\par
\medskip\hrule\medskip\noindent\textbf{Author text begins below.}\par
\setcounter{section}{1}
\setcounter{subsection}{0}
\setcounter{subsubsection}{0}
\setcounter{Theorem}{0}
\setcounter{equation}{0}
% Original source lines 169--850
\section{Definition of the Yangian}\label{sec:1}

Throughout this article we will use the following general conventions.
Let $\mathrm{A}$ and $\mathrm{B}$ be any two associative
$\ZZ_{2}$-graded algebras. Their tensor product
$\mathrm{A}\ot\mathrm{B}$ is also
an associative $\ZZ_2$-graded algebra such that for any homogeneous
elements
$X,X^\prime\in\mathrm{A}$ and $Y,Y^\prime\in\mathrm{B}$
\begin{gather}\label{conv1}
(X\ot Y) (X^\prime\ot Y^\prime)=X X^\prime\ot Y Y^\prime (-1)^{\deg X^\prime\deg Y},\\
\label{conv2}
\deg (X\ot Y)=\deg X+\deg Y .
\end{gather}
For any $\ZZ_2$-graded modules
$U$ and $V$ over $\mathrm{A}$ and $\mathrm{B}$ respectively,
the vector space $U\ot V$ is a~$\ZZ_2$-graded module over
$\mathrm{A}\ot\mathrm{B}$ such that for any homogeneous elements
$x\in U$ and $y\in V$
\begin{gather}\label{XY}
(X\ot Y) ( x\ot y ) =X x\ot Y y (-1)^{\deg x \deg Y},\\
\label{xy}
\deg ( x\ot y )=\deg x+\deg y .
\end{gather}
A homomorphism $\al\colon \mathrm{A}\to\mathrm{B}$ is a linear map
such that $\al (X X^\prime)=\al (X) \al (X^\prime)$ for all
$X,X^\prime\in\mathrm{A} $. But an antihomomorphism
$\be\colon \mathrm{A}\to\mathrm{B}$ is a linear map
such that for all homogeneous $X,X^\prime\in\mathrm{A}$
\begin{gather}\label{bb}
\be (X X^\prime )=\be (X^\prime) \be (X) (-1)^{\deg X\deg X^\prime}.
\end{gather}

If $\mathrm{A}$ is unital,
let $\iota_{h}$ be
its embedding into the tensor product $\mathrm{A}^{\ns\ot n}$
as the $h $-th tensor factor:
\begin{gather*}
\iota_{h}(X)=1^{\ot (h-1)}\ot X\ot1^{\ot (n-h)}\qquad\text{for}\quad h=1,\dots, n .
\end{gather*}
Here $n$ can be any positive integer.
We will also use various embeddings of the algebra
$\mathrm{A}^{\ns\ot m}$ into~$\mathrm{A}^{\ns\ot n}$
for any $m=1,\dots, n$.
For any choice of pairwise distinct indices $h_1,\dots, h_m\in\{1,\dots, n \}$
and of an element $X\in\mathrm{A}^{\ns\ot m}$ of the form
$X=X^{(1)}\ot\cdots\ot X^{(m)}$ we will denote
\begin{gather*}
X_{h_1\dots h_m}=\iota_{h_1}\big(X^{(1)}\big) \cdots \iota_{h_m}\big(X^{(m)}\big) \in\mathrm{A}^{\ns\ot n}.
\end{gather*}
We will then extend the notation $X_{h_1\dots h_m}$ to all elements
$X\in\mathrm{A}^{\ns\ot m}$ by linearity.

Now let the indices $i$, $j$ run through $1,\dots, M+N $.
We will always write $\bi=0$ if $1\le i\le M$ and
$\bi=1$ if $M<i\le M+N $. Consider the
$\ZZ_{2}$-graded vector space $\CMN $.
Let $e_i\in\CMN$ be an element of the standard basis.
The $\ZZ_{2}$-grading on $\CMN$ is defined so that $\deg e_i=\bi $.
Let $E_{ij}\in\End\CMN$ be the standard matrix unit, so that
$E_{ij} e_k=\de_{jk} e_i $. The
associative algebra $\End\CMN$
is $\ZZ_{2}$-graded so that $\deg E_{ij}=\bi+ \bj $.

For any $n$ we can identify
the tensor product $\big(\!\End\CMN\big)^{\ot n}$ with the algebra
$\End\!\big(\big(\CMN\big)^{\!\ot n}\big)$ acting on the vector space $\big(\CMN\big)^{\!\ot n}$
by repeatedly using the conventions~\eqref{XY} and~\eqref{xy}.

Let us introduce the {\it Yangian} of the Lie superalgebra
$\glMN$. This is the complex associative
unital $\ZZ_{2}$-graded algebra $\YMN$ with the countable set of generators
\begin{gather}\label{Tijr}
T^{(r)}_{ij},
\qquad\text{where}\quad r=1,2, \dots \quad\text{and}\quad i ,j=1,\dots, M+N .
\end{gather}
The $\ZZ_{2}$-grading on the algebra $\YMN$
is determined by setting $\deg T_{ij}^{(r)}=\bi+\bj$ for $r\ge1$.
To write down defining relations for these
generators we will employ the series
\begin{gather}\label{3.0}
T_{ij}(u)=\de_{ij}\cdot1+T_{ij}^{(1)} u^{-1}+T_{ij}^{(2)} u^{-2}+\cdots
\end{gather}
in a formal variable $u$ with coefficients from $\YMN $.
Then for all possible indices $i$, $j$, $k$, $l${\samepage
\begin{gather}\label{3.1}
(u-v)[ T_{ij}(u) ,T_{kl}(v) ](-1)^{\bi \bk + \bi \bl + \bk \bl}=T_{kj}(u) T_{il}(v)-T_{kj}(v) T_{il}(u),
\end{gather}
where $v$ is another formal variable.
The square brackets here stand for the supercommutator.
Notice that the series denoted by $T_{ij}(u)$
here and in \cite{N1} differ by the scalar factor
$(-1)^{( \bi+1 ) \bj}$.}

We will also use the following matrix form of the defining relations~\eqref{3.1}.
Take the element
\begin{gather}\label{pop}
P = \sum_{i,j=1}^{M+N} E_{ij}\ot E_{ji} (-1)^{\bj}
\in\big(\End\CMN\big)^{\ot 2} .
\end{gather}
This element acts on the vector space $\big(\CMN\big)^{\ot 2}$ so that
$e_i\ot e_{j}\mapsto e_j\ot e_i {(-1)}^{\bi \bj}$.
Here we identify the algebra $\big(\End\CMN\big)^{\ot 2}$ with the algebra
$\End\big(\big(\CMN\big)^{\ot 2}\big)$ by using~\eqref{XY}.

For any $n$ let $\Sym_n$ be the symmetric group acting
on the set $\{1,\dots, n\}$ by permutations. For each $m=1,\dots, n-1$
denote by $\si_m$ the element of $\Sym_n$ exchanging~$m$ and $m+1 $.
The group $\Sym_n$ also acts on the vector space $\big(\CMN\big)^{\ot n}$.
This action is defined by the assignment
$\si_m\mapsto P_{m,m+1}$ for each~$m $.
Here we identify the algebra
$\big(\End\CMN\big)^{\ot n}$ with $\End\big(\big(\CMN\big)^{\ot n}\big)$ via~\eqref{XY},~\eqref{xy}.

The rational function $R(u)=1-P u\1$ with values in the algebra
$\big(\End\CMN\big)^{\ot 2}$ is called the {\it Yang $R$-matrix}.
It satisfies the {\it Yang--Baxter equation}
in the algebra $\big(\End\CMN\big)^{\ot 3}(u,v,w)$
\begin{gather}\label{3.6}
R_{12}(u-v) R_{13}(u-w) R_{23}(v-w)=R_{23}(v-w) R_{13}(u-w) R_{12}(u-v) .
\end{gather}
Since $P^2=1 $, we also have the relation
\begin{gather}\label{3.52}
R(-u) R(u)=1-u^{-2} .
\end{gather}

Now combine all the series \eqref{3.0} into the single element
\begin{gather}\label{tu}
T(u)=\sum_{i,j=1}^{M+N} E_{ij}\ot T_{ij}(u)
\in\big(\End\CMN\big)\ot\YMN \big[\big[u^{-1}\big]\big] .
\end{gather}
For any $n$ and any $p=1,\dots, n$ we will denote
\begin{gather}
T_p(u)=\iota_p\ot\id ( T(u)) \in\big(\End\CMN\big)^{\ot n}\ot\YMN \big[\big[u\1\big]\big] .\label{3.22}
\end{gather}
By using this notation for $n=2$
the relations \eqref{3.1} can be rewritten as
\begin{gather}\label{3.3}
(R (u-v)\ot1) T_1(u) T_2(v)=T_2(v) T_1(u) (R (u-v)\ot1) .
\end{gather}
Namely, after multiplying each side of \eqref{3.3} by $u-v$ it becomes
a relation of series in $u$, $v$ with coefficients in %the algebra
$\big(\End\CMN\big)^{\ot 2}\ot\YMN$
equivalent to the collection of all the relations~\eqref{3.1}.

\begin{Proposition}\label{m}
An antiautomorphism of $\YMN$ can be defined by the assignment
\begin{gather}\label{M}
T_{ij}(u)\mapsto T_{ij}(-u) .
\end{gather}
\end{Proposition}

\begin{proof}Due to the convention \eqref{conv1}, by using
the notation \eqref{3.22} for $n=2$ we get
\begin{gather}\label{TT}
T_1(u) T_2(v)=\sum_{i,j,k,l=1}^{M+N}
E_{ij}\ot E_{kl}\ot T_{ij}(u) T_{kl}(v)
{(-1)}^{( \bi+ \bj )( \bk+ \bl )} ,\\
\label{MM}
T_2(-v) T_1(-u)=\sum_{i,j,k,l=1}^{M+N}
E_{ij}\ot E_{kl}\ot T_{kl}(-v) T_{ij}(-u) .
\end{gather}
By using the convention \eqref{bb}, the antihomomorphism property of \eqref{M} follows from the relation
\begin{gather*}
(R (u-v)\ot1) T_2(-v) T_1(-u)=
T_1(-u) T_2(-v) (R (u-v)\ot1),
\end{gather*}
which is obtained from~\eqref{3.3} by using~\eqref{3.52}.
The antihomomorphism~\eqref{M} is clearly involutive and therefore bijective.
\end{proof}

For all indices $i$, $j$ define the series $\TP_{ij}(u)$
by using the element inverse to \eqref{tu} so that
\begin{gather}\label{tunat}
T(u)^{-1}=\sum_{i,j=1}^{M+N} E_{ij}\ot\TP_{ij}(u) .
\end{gather}

\begin{Proposition}\label{s}
An antiautomorphism of $\YMN$ can be defined by the assignment
\begin{gather}\label{S}
T_{ij}(u)\mapsto\TP_{ij}(u) .
\end{gather}
\end{Proposition}

\begin{proof}Similarly to \eqref{MM}, by using
the notation \eqref{3.22} for $n=2$ we get
\begin{gather*}
T_2(v)^{-1} T_1(u)^{-1}=\sum_{i,j,k,l=1}^{M+N}E_{ij}\ot E_{kl}\ot\TP_{kl}(v) \TP_{ij}(u) .
\end{gather*}
Comparing this with \eqref{TT} the antihomomorphism property of \eqref{S} follows from the relation
\begin{gather}\label{3.33}
(R (u-v)\ot1) T_2(v)^{-1} T_1(u)^{-1}=
T_1(u)^{-1} T_2(v)^{-1} (R (u-v)\ot1),
\end{gather}
which is obtained by multiplying both sides of the defining relation \eqref{3.3}
on the left and right by~$T_2(v)^{-1}$ and then by~$T_1(u)^{-1}$.
The bijectivity of~\eqref{S} follows from Proposition~\ref{ss} below.
\end{proof}

Further, let $\tau$ be the antiautomorphism of $\End\CMN$
defined by the assignment
\begin{gather*}
E_{ij}\mapsto E_{ji} {(-1)}^{\bi ( \bj + 1)} .
\end{gather*}
Then by the definition \eqref{tu} we have
\begin{gather*}
\tau \ot\id ( T(u) )=
\sum_{i,j=1}^{M+N} E_{ji}\ot T_{ij}(u)
{(-1)}^{\bi ( \bj + 1)}=
\sum_{i,j=1}^{M+N} E_{ij}\ot T_{ji}(u)
{(-1)}^{\bj ( \bi + 1)} .
\end{gather*}

\begin{Proposition}
An antiautomorphism of %the algebra
$\YMN$ can be defined by the assignment
\begin{gather}
\label{T}
T_{ij}(u)\mapsto T_{ji}(u)
{(-1)}^{\bj ( \bi + 1)} .
\end{gather}
\end{Proposition}

\begin{proof}
Observe that $(\tau\ot\tau)(P)=P$ and hence
\begin{gather}
\label{Rtt}
(\tau\ot\tau)(R(u-v))=R(u-v) .
\end{gather}
Therefore the antihomomorphism property of \eqref{T}
follows from the relation which is obtained by applying
$\tau\ot\tau\ot\id$ to both sides of\eqref{3.3},
see the proofs of Propositions \ref{m} and \ref{s} above.
To prove the bijectivity of
the antihomomorphism \eqref{T}, observe
that its square is given by %the assignment
\begin{gather*}
T_{ij}(u)\mapsto T_{ij}(u)
{(-1)}^{\bi+ \bj} .
\end{gather*}
In particular, the square is an automorphism of the $\ZZ_2$-graded
algebra $\YMN $.
\end{proof}

Put
$\TS(u)=\tau \ot\id \big( T(u)^{-1} \big) $.
Then by \eqref{tunat}
\begin{gather}\label{tsu}
\TS(u)=
\sum_{i,j=1}^{M+N} E_{ji}\ot\TP_{ij}(u)
{(-1)}^{\bi ( \bj + 1)}=
\sum_{i,j=1}^{M+N} E_{ij}\ot\TP_{ji}(u)
{(-1)}^{\bj ( \bi + 1)} .
\end{gather}

\begin{Corollary}
\label{ts}
An automorphism of the algebra $\YMN$ can be defined by the assignment
\begin{gather}
\label{tusharp}
T_{ij}(u)\mapsto\TP_{ji}(u)
{(-1)}^{\bj ( \bi + 1)} .
\end{gather}
\end{Corollary}

\begin{proof}
The assignment \eqref{tusharp} can also be obtained by first applying
\eqref{T} to $T_{ij}(u)$ and then applying \eqref{S} to the result.
Hence \eqref{tusharp} defines an automorphism of the algebra
$\YMN$ as a composition of two antiautomorphisms.
\end{proof}

By the definition \eqref{tsu} the homomorphism property of %the assignment
\eqref{tusharp} is equivalent to the relation
\begin{gather}
\label{3.333}
(R (u-v)\ot1) \TS_1(u) \TS_2(v)=
\TS_2(v) \TS_1(u)(R (u-v)\ot1),
\end{gather}
which can also be obtained by applying
$\tau\ot\tau\ot\id$ to both sides of \eqref{3.33}
and then using~\eqref{Rtt}.

\begin{Proposition}\label{mst}
The antiautomorphisms \eqref{M}, \eqref{S}, \eqref{T} of~$\YMN$ pairwise commute.
\end{Proposition}

\begin{proof}The antiautomorphism \eqref{M} clearly commutes with either of~\eqref{S},~\eqref{T}.
To prove the commutativity of the latter two, consider the tensor product of the antiautomorphisms~$\tau\1$ and~\eqref{T}
of $\End\CMN$ and $\YMN$ respectively. This is
is an antiautomorphism of the
algebra $\End\CMN\ot\YMN $, and the series \eqref{tu} is invariant
under this antiautomorphism. Hence the series \eqref{tunat}
is also invariant under it.
The latter invariance implies that~\eqref{T} maps
\begin{gather*}%\label{ST}
\TP_{ij}(u)\mapsto\TP_{ji}(u) {(-1)}^{\bj ( \bi + 1)} .
\end{gather*}
Hence \eqref{tusharp} can also be obtained by first applying~\eqref{S} to $T_{ij}(u)$ and then applying~\eqref{T} to the result.
Comparing this with the proof of Corollary~\ref{ts} completes the argument.
\end{proof}

Consider the universal enveloping
algebra $\UMN$ of the Lie superalgebra $\glMN $.
To avoid confusion, the element of $\glMN$
corresponding to $E_{ij}\in\End\CMN$ will be denoted by $e_{ij} $.
By definition, the bracket on $\glMN$ is the supercommutator.
Hence in~$\UMN$
\begin{gather}\label{glMN}
[ e_{ij} ,e_{kl} ]=
\de_{jk} e_{il}-
\de_{li} e_{kj}
{(-1)}^{( \bi+ \bj )( \bk+ \bl )} .
\end{gather}
By using the defining relations \eqref{3.1} one can demonstrate that
there is a homomorphism
\begin{gather}\label{3.5}
\YMN\to\UMN\colon \ T_{ij}(u)\mapsto \de_{ij}-e_{ji} u^{-1} {(-1)}^{\bj} .
\end{gather}
This homomorphism is surjective. The relations \eqref{3.1} also imply that there is a homomorphism
\begin{gather}\label{3.51}
\UMN\to\YMN\colon \
e_{ji}\mapsto- T_{ij}^{(1)} {(-1)}^{\bj} .
\end{gather}
The composition of homomorphisms \eqref{3.51}
and~\eqref{3.5}
is the identity map $\UMN\to\UMN $. So
\eqref{3.51} is an embedding of $\ZZ_2$-graded associative unital algebras.
The homomorphism~\eqref{3.5} is identical on the subalgebra~$\UMN $.
It is called
the {\it evaluation homomorphism} for $\YMN $.

There is a natural Hopf algebra structure on $\YMN $.
A coassociative comultiplication homomorphism $\De\colon \YMN\to\YMN\ot\YMN$
can be defined by the assignment
\begin{gather}\label{3.7}
T_{ij}(u)\mapsto\sum_{k=1}^{M+N}
T_{ik}(u)\ot T_{kj}(u)
{(-1)}^{( \bi + \bk )( \bj + \bk )},
\end{gather}
where the tensor product is taken over the subalgebra $\CC\big[\big[u^{-1}\big]\big]$
in $\YMN \big[\big[u\1\big]\big] $.
The counit homomorphism $\ep\colon \YMN\to\CC$ is defined
by the assignment $T_{ij}(u)\mapsto\de_{ij} $.
The antipodal mapping $\S\colon \YMN\to\YMN$ is the antiautomorphism~\eqref{S}.
Justification of all these
definitions is similar to that in the case $N=0$
considered for instance in \cite[Section~1]{MNO}. Here we omit the details.
Note that \eqref{3.51} is an embedding of Hopf algebras as
by the above definitions
\begin{gather*}
\De\colon \ T_{ij}^{(1)}\mapsto T_{ij}^{(1)}\ot1+1\ot T_{ij}^{(1)} ,
\qquad
\ep\colon \ T_{ij}^{(1)}\mapsto0 ,
\qquad
\S\colon \ T_{ij}^{(1)}\mapsto- T_{ij}^{(1)} .
\end{gather*}

There are two natural
ascending filtrations on
the associative algebra $\YMN $. The first one is defined by
assigning the degree~$r$ to the generator~\eqref{Tijr}.
Let $\gr\YMN$ be the corresponding graded algebra.

Let us denote by $X_{ij}^{(r)}$ the image of $T_{ij}^{(r)}$
in the degree $r$ component of $\gr\YMN $.
Observe that the $\ZZ_2 $-grading on the algebra $\YMN$ descends to
$\gr\YMN$ so that the degree of the image is again
$\bi + \bj $. It follows from the relations \eqref{3.1}
that these images supercommute.
We shall prove that
$\gr\YMN$ is a free supercommutative algebra generated by
these images.

Now introduce another filtration on the associative algebra $\YMN$
by assigning the deg\-ree~$r-1$ to the generator~\eqref{Tijr}.
Let $\grpr\YMN$ be the corresponding graded algebra.
Consider the latter algebra.

Let us denote by $Y_{ij}^{(r)}$ the image of $T_{ij}^{(r)}$
in the degree $r-1$ component of $\grpr\YMN $.
The $\ZZ_2 $-grading on the algebra
$\YMN$ descends to
$\grpr\YMN$ so that the degree of the image is~$\bi + \bj $.
The graded algebra $\grpr\YMN$ inherits from $\YMN$ the Hopf algebra structure too.
Namely, by the above definitions in
the graded Hopf algebra $\grpr\YMN$ for $r\ge1$ %we get
\begin{gather}
\label{grprH}
\De: Y_{ij}^{(r)}\mapsto Y_{ij}^{(r)}\ot1+1\ot Y_{ij}^{(r)} ,
\ \quad
\ep: Y_{ij}^{(r)}\mapsto0 ,
\ \quad
\S: Y_{ij}^{(r)}\mapsto- Y_{ij}^{(r)} .
\end{gather}

Let us now consider the polynomial current Lie superalgebra $\g $.
The elements $e_{ji} u^{r}$ with $r=0 ,1 ,2 , \dots$
and $i ,j=1 ,\dots, M+N$ make a basis of~$\g $.
The $\ZZ_{2}$-grading on~$\g$ is defined by
$\deg e_{ji} u^{r}=\bi+\bj $. Take
the universal enveloping algebra $\Ug $.

\begin{Proposition}
\label{P21}
One can define a surjective homomorphism $\Ug\to\grpr\YMN$
of $\ZZ_2$-graded associative algebras by mapping for $r\ge0$
\begin{gather}
\label{UY}
e_{ji} u^r\mapsto- Y_{ij}^{(r+1)} {(-1)}^{\bj} .
\end{gather}
\end{Proposition}

\begin{proof}
Due to \eqref{glMN} the supercommutation relations with $r ,s\ge0$
\begin{gather}
\label{grel}
[ e_{ji} u^{r}, e_{lk} u^{s} ]=
\de_{il} e_{jk} u^{r+s}-
\de_{kj} e_{li} u^{r+s}
{(-1)}^{( \bi+ \bj )( \bk+ \bl )}
\end{gather}
define the Lie superalgebra $\g $.
Multiplying the relation \eqref{grel} by
$(-1)^{\bi \bk + \bi \bl + \bk \bl+\bj+\bl}$
and replacing the basis elements of $\g$ there by their images
under the mapping \eqref{UY} we get
\begin{gather*}
[ Y_{ij}^{(r+1)},Y_{kl}^{(s+1)} \bigr]
(-1)^{\bi \bk + \bi \bl + \bk \bl}=
\de_{kj} Y_{il}^{(r+s+1)}-\de_{il} Y_{kj}^{(r+s+1)} .
\end{gather*}
The latter relation in $\grpr\YMN$ follows from \eqref{3.1},
see for instance \cite[Section 1]{MNO}.
This proves the homomorphism property of the assignment \eqref{UY}.
This homomorphism is clearly surjective
and preserves the $\ZZ_{2}$-grading.
\end{proof}

It immediately follows from \eqref{grprH} that
\eqref{UY} is a homomorphism of Hopf algebras.
Here we use the standard Hopf algebra structure on $\Ug$
as the universal enveloping algebra of a Lie superalgebra.
We shall demonstrate that the homomorphism \eqref{UY} is also injective.

By comparing \eqref{3.6} with \eqref{3.3}
we obtain that for any $z\in\CC$ the assignment
\begin{gather}\label{3.66}
\End\CMN\ot\YMN \big[\big[u^{-1}\big]\big]
\to
\big(\End\CMN\big)^{\ot 2} \big[\big[u^{-1}\big]\big]\colon \
T(u)\mapsto R(u-z)
\end{gather}
defines a representation $\YMN\to\End\CMN$. More explicitly, under
\eqref{3.66} for any $r\ge0$
\begin{gather}
\label{2.19}
T_{ij}^{(r+1)}\mapsto
- E_{ji} z^{r} {(-1)}^{\bj} .
\end{gather}
Note that this representation of $\YMN$
can be also obtained from the standard representation $\UMN\to\End\CMN$
by pulling back through the evaluation homomorphism \eqref{3.5}
and then back through the automorphism of $\YMN$ defined by mapping
$
T_{ij}(u)\mapsto T_{ij}(u-z) .
$

The comultiplication on $\YMN$
now allows us to define for $n=1 ,2 , \dots$ a representation
$\YMN\to(\End\CMN)^{\ot n}$ depending on
$z_1,\dots, z_n\in\CC $. This is the tensor product of
the representations \eqref{3.66} where $z=z_1,\dots, z_n $.
Due to \eqref{tu} and \eqref{3.7} under this representation
\begin{gather}\label{3.666}
T(u)\mapsto R_{12}(u-z_1) \cdots R_{1,n+1}(u-z_n) .
\end{gather}

\begin{Proposition}\label{P22}
Let the complex parameters $z_1,\dots, z_n$ and the positive integer~$n$ vary.
Then the kernels of the representations~\eqref{3.666} of~$\YMN$
have the zero intersection.
\end{Proposition}

\begin{proof}Take any finite linear combination of the products
$T_{i_1j_1}^{(r_1+1)}\cdots T_{i_mj_m}^{(r_m+1)}\in\YMN$ with
$
A_{i_1j_1\dots i_mj_m}^{r_1\dots r_m}\in\CC
$
being the respective coefficients.
In this linear combination the number $m\ge0$ and
indices $r_1,\dots, r_m\ge0$ may vary.
Consider the image of this linear combination
under the representation
$\YMN\to\big(\End\CMN\big)^{\ot n}$ defined by \eqref{3.666}.
This image depends on $z_1,\dots, z_n$ polynomially.
Let $A$ be the sum of those terms of
this polynomial which have the maximal total degree in $z_1,\dots, z_n $.
Let $d$ be this degree.

Consider the second of our two ascending filtrations
on the associative algebra $\YMN $,
the corresponding graded algebra being $\grpr\YMN $.
Equip the tensor product $\YMN^{\ot n}$ with the ascending
filtration where the degree is the sum of the degrees on the tensor factors.
Then by the definition \eqref{3.7}
under the comultiplication $\YMN\to\YMN^{\ot n}$ for $r\ge0$
\begin{gather*}
T_{ij}^{(r+1)}\mapsto
\sum_{h=1}^n
1^{\ot (h-1)}\ot T_{ij}^{(r+1)}\ot1^{\ot (n-h)}
 +
\text{lower degree terms.}
\end{gather*}
Therefore the sum $A\in\big(\End\CMN\big)^{\ot n}$
coincides with the image of the sum
\begin{gather*}
\sum_{r_1+\dots+r_m= d}
A_{i_1j_1\dots i_mj_m}^{r_1\dots r_m}
e_{j_1i_1} u^{r_1}\cdots e_{j_mi_m} u^{r_m}
(-1)^{m + \bj_1 + \dots + \bj_m}\in\Ug
\end{gather*}
under the tensor product of the evaluation representations
\begin{gather}\label{evalrep}
\Ug\to\End\CMN\colon \
e_{ji} u^{r}
\mapsto
E_{ji} z^{r}
\end{gather}
at the points $z=z_1,\dots, z_n $. Here we used the explicit
description~\eqref{2.19}
of the representation $\YMN\to\End\CMN$ corresponding to $z\in\CC $.
Due to Proposition~\ref{P21} it now
suffices to show that when the complex parameters $z_1,\dots, z_n$
and the positive integer~$n$ vary, the kernels
of the tensor products
of the evaluation representations
of the algebra $\Ug$ at $z=z_1,\dots, z_n$
have the zero intersection.
This will also imply that the homomorphism~\eqref{UY} is injective.

Choose any $\ZZ_2 $-homogeneous
basis $f_1,\dots, f_{(M+N)^2}$ of $\glMN$
such that the first vector $f_1$ is
\begin{gather}\label{esum}
\sum_{i=1}^{M+N}e_{ii} .
\end{gather}
The corresponding elements of $\End\CMN$
will be denoted by $F_1,\dots, F_{(M+N)^2} $.
Hence $F_1=1 $.
The elements $f_{p} u^{r}$
with $r=0 ,1 ,2 , \dots$
and $p=1 ,\dots, (M+N)^2$ constitute a basis of $\g $.
Choose any total ordering of this basis which ends with the
infinite sequence
$\dots , f_1 u^{2}, f_1 u , f_1$.
Take any finite linear combination $L$ of the products{\samepage
\begin{gather}\label{2.3.1}
f_{p_1}u^{r_1}\cdots f_{p_m}u^{r_m}\in\Ug
\end{gather}
with $L_{p_1\dots p_m}^{r_1\dots r_m}\in\CC$
being the coefficients.
We assume that the factors in the products are
arranged according to our ordering of the basis of $\g $.
Due to the supercommutation relations in $\Ug$
we assume it without any loss of generality.
The basis elements of $\ZZ_2 $-degree~$0$
may occur in any product~\eqref{2.3.1} with a multiplicity
but those of $\ZZ_2 $-degree $1$ may occur at most~once.

}

Let us denote by $\rho_z$ the evaluation representation~\eqref{evalrep}.
More generally, denote by $\rho_{z_1\dots z_n}$ the tensor product
$\rho_{z_1}\ot\cdots\ot\rho_{z_n}$
pulled back through $n $-fold comultiplication on $\Ug $.
Hence $\rho_{z_1\dots z_n}$ is a homomorphism of
associative algebras
\begin{gather*}
\Ug\to\big(\End\CMN\big)^{\ot n}\colon \
f_{p} u^{r}\mapsto
F_{p} z_1^{r}\ot1^{\ot (n-1)}+\dots+
1^{\ot (n-1)}\ot F_{p} z_n^{r} .
\end{gather*}
Suppose that $\rho_{z_1\dots z_n}(L)=0$ for every $n$
and all $z_1,\dots, z_n\in\CC $.
We need to prove that $L=0 $.

For each product \eqref{2.3.1} there is a number $a$ such that
$p_1,\dots, p_{a}>1$ but $p_{a+1} ,\dots, p_m=1 $.
This is due to our ordering of the basis of $\g $.
The numbers $a$ for different products
\eqref{2.3.1} may differ, and we do not exclude the case $a=0$.
Let $h$ be the maximum of the numbers $a $.

Suppose $n\ge h $.
Let $\om_{h}$ be the supersymmetrisation map of
the tensor product $\glMN[u]^{\ot h}$ normalised so that
$\om_{h}^{2}=h ! \om_{h} $.
Let $W$ be the subspace
of $\big(\End\CMN\big)^{\ot n}$ spanned by the vectors
$F_{q_{1}}\ot\cdots\ot F_{q_{n}}$
where at least one
of the indices $q_{1},\dots, q_{h}$ is $1 $.
If $h=0$ then this subspace is assumed to be zero.
Applying the homomorphism
$\rho_{z_1\dots z_n}$ to a product \eqref{2.3.1}
with $a=h$ gives
\begin{gather}\label{2.3.2}
( \rho_{z_1}\ot\cdots\ot\rho_{z_{h}} )
\big( \om_{h}
\big( f_{p_1}u^{r_1}
\ot\cdots\ot
f_{p_{h}}u^{r_h} \big) \big)
\ot1^{\ot (n-h)}
\prod_{b=h+1}^m
\big( z_1^{r_b}+\dots+z_n^{r_b} \big)
\end{gather}
modulo the subspace $W$. Applying
$\rho_{z_1\dots z_n}$ to a product~\eqref{2.3.1}
with $a<h$ gives an element~of~$W$.
But a linear combination of the
expressions~\eqref{2.3.2} belongs to~$W$
only if this combination is zero.

For each product \eqref{2.3.1} with $a=h$
there is a number $c\ge h$ such that $r_{h+1}\ge\cdots\ge r_c>0$
but $r_{c+1} ,\dots, r_m=0 $.
This is due to our ordering of the basis of~$\g $.
Then~\eqref{2.3.2} equals
\begin{gather}\label{2.3.3}
( \rho_{z_1}\ot\dots\ot\rho_{z_{h}} )
\big( \om_{h}
\big( f_{p_1}u^{r_1}
\ot\dots\ot
f_{p_{h}}u^{r_h} \big) \big)
\ot1^{\ot (n-h)}
\prod_{b=h+1}^c
\big( z_1^{r_b}+\dots+z_n^{r_b} \big)
\end{gather}
multiplied by $n^{m-c}$.
Let $g$ be the maximum of the numbers $c$
for all products~\eqref{2.3.1} with $a=h $.

Suppose $n\ge g$. Consider the pairs of sequences of indices
$p_1,\dots, p_{m}$ and $r_1,\dots, r_m$
showing in $L$ in the products \eqref{2.3.1}
with $a=h$. For every such a pair there is
$c\in\{h,\dots, g \}$.
Then we have
$p_{h+1},\dots, p_{m}=1$ and $r_{c+1},\dots, r_m=0$.
Take all different pairs of sequences $p_1,\dots, p_{h}$ and $r_1,\dots, r_c$
arising in this way.
The expressions~\eqref{2.3.3} corresponding to the latter pairs
are linearly independent as polynomials in $z_1,\dots, z_n$
with values in $\big(\End\CMN\big)^{\ot n}$.
This is again due to our ordering of the basis of~$\g$.
Here we also employ the observation that~in~\eqref{2.3.3} the image of
$\rho_{z_1}\ot\cdots\ot\rho_{z_{h}}$
does not depend on the parameters $z_{h+1},\dots, z_n$
while the product over $b=h+1,\dots, c$ in~\eqref{2.3.3}
does depend on these parameters whenever $c>h$.
Therefore if $\rho_{z_1\ldots z_n}(L)=0$ for a certain $n\ge g$
and for all $z_1,\dots, z_n\in\CC$ then for each of the latter pairs
\begin{gather*}
\sum_{m=c}^\infty
L_{p_1\ldots p_h 1\ldots 1}^{r_1\ldots r_c 0\ldots 0} n^{m-c}=0.
\end{gather*}
There are exactly $m$ lower indices and also $m$ upper indices
in the coefficient
$L_{p_1\ldots p_h 1\ldots 1}^{r_1\ldots r_c 0\ldots 0}$ above.
By letting the number $n$ vary we now prove that all these
coefficients vanish.
\end{proof}

In the course of the proof of Proposition \ref{P22} we established that
the homomorphism~\eqref{UY} is injective.
Together with Proposition~\ref{P21}
and the observation made just after its proof this~yields

\begin{Theorem}\label{T18}The Hopf algebras $\Ug$ and $\grpr\YMN$ are isomorphic via~\eqref{UY}.
\end{Theorem}
\begin{thebibliography}{99}
\bibitem[6]{MNO}
Molev A., Nazarov M., Olshanski G., Yangians and classical {L}ie algebras,
 \href{https://doi.org/10.1070/RM1996v051n02ABEH002772}{\textit{Russian Math. Surveys}} \textbf{51} (1996), 205--282,
 \href{https://arxiv.org/abs/hep-th/9409025}{arXiv:hep-th/9409025}.

\bibitem[8]{N1}
Nazarov M.L., Quantum {B}erezinian and the classical {C}apelli identity,
 \href{https://doi.org/10.1007/BF00401646}{\textit{Lett. Math. Phys.}} \textbf{21} (1991), 123--131.

\end{thebibliography}
\end{document}
